
\chapterstar{Planteamiento del Problema}

Muchas pequeñas y medianas empresas dependen de hojas de cálculo o de sistemas incompletos para gestionar inventarios y ventas, lo que genera errores frecuentes, duplicidad de información y dificultades para obtener datos en tiempo real. 
La falta de un control centralizado entre el inventario físico y el registrado provoca pérdidas económicas por productos caducados, desabasto, por falta de existencias o exceso de inventario que no se coloca a tiempo, esto también crea problemas de concurrencia, ya que es difícil saber quien, cuando y porque se modifico un dato. La falta de integración tecnológica no solo limitan su actividad diaria, sino que también representa una amenaza en la competitividad en un mercado cada vez mas tecnológico, el núcleo del problema radica en la ausencia de un sistema de información unificado y en tiempo real.
En los negocios que manejan productos sensibles a condiciones ambientales, como los farmacéuticos o alimenticios, la ausencia de un sistema que integre el monitoreo físico de variables de temperatura y/o humedad puede complicar la situación de estos mismos. En el sector farmacéutico o alimenticio se debe llevar un control de temperatura y caducidad, la falta de un registro continuo de estas condiciones bloquearía la capacidad de prever un error, a su vez, corregir dicho error.
Las soluciones de gestión de mercado para las pequeñas y medianas empresas a menudo son caras o carecen de funcionalidades modernas esenciales, en un mundo post-pandemia, la gestión requiere acceso la mayoría sino es que todo el tiempo desde cualquier dispositivo, los sistemas locales y de escritorio limitan la capacidad de gestión del administrador cuando esta fuera de la empresa.
Las plataformas actuales no suelen ofrecer una automatización de alertas sofisticada que combine múltiples factores como falta de inventario, temperatura alta o baja en los diferentes puntos de la empresa, la gestión se mantiene con el sistema manual, requiriendo supervisión constante.
\thispagestyle{empty}
